% Eigenständiger Lesetext zum Einbettungssatz für kommutative kürzbare Monoide.
% Die konkrete Ganzzahlkonstruktion wird in der Lesefassung zu Band 17 entwickelt.
\par\Needspace{10\baselineskip}
\hypertarget{differences.reading.problem}{}
\section*{Vom Zahlenbeispiel zur allgemeinen Frage}

Die Gleichung \(x+5=2\) hat in den natürlichen Zahlen keine Lösung.
In den ganzen Zahlen hat sie die Lösung \(x=-3\). Beim Übergang
von \(\mathbb N\) zu \(\mathbb Z\) wird die bisherige Addition
erhalten, und jedes Element erhält ein additives Gegenstück.
Dadurch werden Gleichungen lösbar, für die der ursprüngliche Bereich
zu klein war.

Hinter diesem vertrauten Schritt steht eine allgemeine Frage:
Welche Rechengesetze erlauben es, eine vorhandene Addition um
Gegenstücke zu erweitern, ohne dabei verschiedene ursprüngliche
Elemente gleichzusetzen? Die Elemente müssen keine Zahlen sein.
Auch Paare natürlicher Zahlen, die man koordinatenweise addiert,
geben ein Beispiel. Der Aufbau der Erweiterung wird in beiden
Fällen von denselben Gesetzen getragen.

Diese Lesefassung entwickelt den allgemeinen Einbettungssatz für
\emph{kommutative kürzbare Monoide}. Wir erklären die Voraussetzungen
und führen die gesamte Konstruktion durch: von Paaren über
Äquivalenzklassen zu einer abelschen Gruppe. Der Hauptsatz bleibt
in Band~41 unter
\FormulaRefAuto{GroupCompletionEmbeddingTheorem}[thm]; die
ausführlichen formalen Ableitungen stehen in den
\DifferencesProofsLink{}.

Die ergänzende \IntegersReadingLink{} zu Band~17 verfolgt den
konkreten Zahlenaufbau weiter. Dort werden neben Addition und
Negation auch Multiplikation, Ordnung und zweiseitige Induktion
entwickelt und die Ganzzahlaxiome im konstruierten Modell
nachgewiesen. Hier liegt der Schwerpunkt auf den allgemeinen
Voraussetzungen und darauf, weshalb die Differenzenkonstruktion
für jede passende Struktur funktioniert. Beide Texte können
unabhängig voneinander gelesen werden.

\par\Needspace{10\baselineskip}
\hypertarget{differences.reading.general}{}
\section*{Die Voraussetzungen und das Ziel}

Sei \(M\) eine Menge mit einer Operation
\(+\colon M\times M\to M\) und einem Element \(0\in M\).
Wir verwenden hier \emph{additive Schreibweise}: Das Zeichen \(+\)
steht für die gegebene Operation, \(0\) für ihr neutrales Element.
Es muss sich dabei nicht um die Addition von Zahlen handeln.
Für alle \(a,b,c\in M\) setzen wir voraus:
\[
  \begin{aligned}
    (a+b)+c&=a+(b+c), &&\text{Assoziativität,}\\
    a+0&=a=0+a, &&\text{neutrales Element,}\\
    a+b&=b+a, &&\text{Kommutativität,}\\
    a+c=b+c&\ \Longrightarrow\ a=b,
      &&\text{Kürzbarkeit.}
  \end{aligned}
\]
Eine Menge mit einer assoziativen Operation und einem neutralen
Element heißt ein \emph{Monoid}. Mit den beiden weiteren Gesetzen
ist \((M,+,0)\) ein \emph{kommutatives kürzbares Monoid}.
Wegen der Kommutativität erlaubt das Kürzungsgesetz auch, einen
gemeinsamen Summanden auf der linken Seite zu kürzen.

Kürzen bedeutet noch nicht, dass ein Element abgezogen werden kann.
Das Gesetz besagt nur: Zwei Summen mit demselben zusätzlichen
Summanden können nur dann gleich sein, wenn die übrigen Summanden
gleich sind. In \(\mathbb N\) gilt dieses Gesetz, obwohl dort
nicht jede Subtraktion möglich ist.

Eine \emph{Gruppe} besitzt zusätzlich zu den Monoidgesetzen zu
jedem Element ein beidseitiges Inverses. In additiver Schreibweise
ist das ein Element, dessen Summe mit dem gegebenen Element null
ergibt. Eine Gruppe mit kommutativer Operation heißt
\emph{abelsch}. Unser Ziel lautet:

\medskip\noindent\textbf{Einbettungssatz.}
Zu jedem kommutativen kürzbaren Monoid \((M,+,0)\) gibt es eine
abelsche Gruppe \((K(M),\oplus,0_K)\) und eine injektive Funktion
\(\iota\colon M\to K(M)\) mit
\[
  \iota(0)=0_K,
  \qquad
  \iota(a+b)=\iota(a)\oplus\iota(b)
  \quad(a,b\in M).
\]

Die beiden Gleichungen sagen, dass \(\iota\) die bestehende
Struktur erhält. Die Injektivität sichert, dass verschiedene
Elemente von \(M\) verschieden bleiben. Eine Abbildung mit diesen
drei Eigenschaften heißt eine \emph{Monoideinbettung}. Sie macht
das ursprüngliche Monoid als Kopie innerhalb der Gruppe verfügbar.
Die folgenden Abschnitte beweisen den Satz durch eine ausdrückliche
Konstruktion von \(K(M)\) und \(\iota\).

\par\Needspace{10\baselineskip}
\hypertarget{differences.reading.pairs}{}
\section*{Formale Differenzen und ihre Gleichheit}

Ein Paar \((a,b)\in M\times M\) soll eine formale Differenz
darstellen: das Ergebnis, das wir erhalten möchten, wenn wir von
\(a\) das Element \(b\) abziehen. Im Monoid ist eine solche
Subtraktion noch nicht gegeben. Das Paar ist dagegen bereits ein
vorhandenes Objekt.

Verschiedene Paare sollen dieselbe Differenz darstellen dürfen.
Wir legen ihre Gleichheit durch eine Bedingung fest, die nur die
gegebene Operation verwendet:
\[
  (a,b)\sim(c,d)
  \quad\Longleftrightarrow\quad
  a+d=b+c.
\]
Die Gleichung heißt \emph{Kreuzsummenbedingung}. Sie drückt die
beabsichtigte Gleichheit von Differenzen aus, ohne eine bereits
existierende Subtraktion vorauszusetzen. Zum Beispiel gilt für
jedes \(t\in M\)
\[
  (a,b)\sim(a+t,b+t),
\]
denn \(a+(b+t)=b+(a+t)\). Dazu verwenden wir Assoziativität
und Kommutativität.

Damit diese Regel Paare zu gleichwertigen Darstellungen
zusammenfassen kann, muss sie eine Äquivalenzrelation sein.
Wir weisen ihre drei Eigenschaften nach.

\medskip\noindent\textbf{Reflexivität und Symmetrie.}
Wegen \(a+b=b+a\) ist jedes Paar zu sich selbst äquivalent.
Aus \(a+d=b+c\) folgt durch Umordnen und Umkehren der Gleichung
\(c+b=d+a\). Somit kann jede Äquivalenz auch in umgekehrter
Richtung gelesen werden.

\medskip\noindent\textbf{Transitivität.}
Es gelte \((a,b)\sim(c,d)\) und \((c,d)\sim(u,v)\), also
\[
  a+d=b+c,
  \qquad c+v=d+u.
\]
Gesucht ist \(a+v=b+u\). Zunächst rechnen wir mit einem
gemeinsamen zusätzlichen Summanden:
\[
  \begin{aligned}
    (a+v)+d
      &=(a+d)+v\\
      &=(b+c)+v\\
      &=b+(c+v)\\
      &=b+(d+u)\\
      &=(b+u)+d.
  \end{aligned}
\]
Die Rechnung verwendet Assoziativität, Kommutativität und die
beiden Voraussetzungen. Jetzt wird die Kürzbarkeit benötigt:
Aus \((a+v)+d=(b+u)+d\) folgt \(a+v=b+u\).
Damit gilt \((a,b)\sim(u,v)\). Wir haben \(d\) nach einem
vorausgesetzten Gesetz gekürzt; ein Inverses von \(d\) wurde
nicht verwendet.

Dies beweist die Äquivalenzrelation aus
\FormulaRefAuto{GroupCompletionPairEquivalence}[thm].

\medskip\noindent\textbf{Äquivalenzklassen als neue Objekte.}
\hypertarget{differences.reading.classes}{}
Wir fassen nun alle zu einem Paar äquivalenten Paare zu einer
Menge zusammen und schreiben
\[
  [a,b]:=\{(c,d)\in M\times M\mid(c,d)\sim(a,b)\}.
\]
Diese Menge heißt die \emph{Äquivalenzklasse} von \((a,b)\).
Das Paar ist ein \emph{Repräsentant} der Klasse. Die Menge aller
Klassen ist unser neuer Bereich:
\[
  K(M):=(M\times M)/{\sim}
       =\{[a,b]\mid a,b\in M\}.
\]
Jede Klasse ist eine Teilmenge von \(M\times M\). Der Quotient
\(K(M)\) ist also eine Menge von Teilmengen dieses Produktraums.

Für seine Elemente gilt die entscheidende Charakterisierung
\[
  [a,b]=[c,d]
  \quad\Longleftrightarrow\quad a+d=b+c.
\]
Sind die Paare äquivalent, so ist jedes zu \((a,b)\) äquivalente
Paar durch Transitivität auch zu \((c,d)\) äquivalent. Symmetrie
liefert die umgekehrte Folgerung. Beide Klassen enthalten daher
genau dieselben Paare. Sind umgekehrt die Klassen gleich, so liegt
\((a,b)\) wegen der Reflexivität in seiner eigenen Klasse und
damit auch in \([c,d]\). Das bedeutet gerade
\((a,b)\sim(c,d)\).
So folgt \FormulaRefAuto{GroupCompletionClassEquality}[thm].

Im besonderen Fall \(M=\mathbb N\) ist diese Konstruktion das
additive Quotientenmodell der ganzen Zahlen. Beispielsweise gilt
\([2,5]=[3,6]=[0,3]\). Jede Gleichheit wird mit natürlichen
Kreuzsummen geprüft. In der folgenden Zeichnung bezeichnet
\([a,b]\) eine Klasse dieses Zahlenbeispiels.

\DifferencesPairGrid

\par\Needspace{10\baselineskip}
\hypertarget{differences.reading.addition}{}
\section*{Eine Operation auf Klassen}

Wir kehren zum allgemeinen Monoid \(M\) zurück. Zwei formale
Differenzen sollen addiert werden, indem wir jeweils ihre ersten
und ihre zweiten Einträge addieren:
\[
  [a,b]\oplus[c,d]:=[a+c,b+d].
\]
Das Zeichen \(\oplus\) unterscheidet die neue Operation auf
Klassen von der gegebenen Operation \(+\) auf \(M\).
Weil \(a+c,b+d\in M\) sind, ist die rechte Seite eine Klasse
in \(K(M)\). Doch bislang haben wir nur einen Ergebnisvorschlag
für gewählte Repräsentanten. Wir müssen zeigen, dass jede andere
Wahl dieselbe Ergebnisklasse liefert.

Nehmen wir deshalb an, dass beide Repräsentanten ausgetauscht werden:
\[
  [a,b]=[a',b'],
  \qquad [c,d]=[c',d'].
\]
Nach der Klassencharakterisierung bedeutet dies
\[
  a+b'=b+a',
  \qquad c+d'=d+c'.
\]
Für die vorgeschlagenen Ergebnispaare rechnen wir:
\[
  \begin{aligned}
    (a+c)+(b'+d')
      &=(a+b')+(c+d')\\
      &=(b+a')+(d+c')\\
      &=(b+d)+(a'+c').
  \end{aligned}
\]
Die erste und die letzte Gleichheit ordnen die Summanden mit
Assoziativität und Kommutativität um. Die mittlere Gleichheit
verwendet die beiden Voraussetzungen. Dies ist die
Kreuzsummenbedingung für
\[
  (a+c,b+d)\sim(a'+c',b'+d').
\]
Also ist \([a+c,b+d]=[a'+c',b'+d']\). Der Ergebniswert
bleibt beim Austausch beider Repräsentanten unverändert.
Das ist der Inhalt von
\FormulaRefAuto{GroupCompletionOperationCompatible}[thm].

Jede Klasse hat nach ihrer Definition mindestens einen
Repräsentanten. Für jedes Paar von Klassen existiert daher ein
Ergebnisvorschlag. Der vorstehende Beweis zeigt, dass alle diese
Vorschläge gleich sind. Die Regel bestimmt somit genau einen Wert
für jedes Klassenpaar und definiert eine Funktion
\[
  \oplus\colon K(M)\times K(M)\longrightarrow K(M).
\]
Dies ist der Übergang von einer Operation auf Darstellungen zu
einer \emph{wohldefinierten} Operation auf den dargestellten
Objekten. Man muss dafür keine bevorzugten Repräsentanten
auswählen. Der formale Existenz- und Eindeutigkeitsnachweis steht
in \FormulaRefAuto{GroupCompletionOperationQuotientDescent}[thm].

\par\Needspace{10\baselineskip}
\hypertarget{differences.reading.group}{}
\section*{Die Klassen bilden eine abelsche Gruppe}

Weil \(\oplus\) wohldefiniert ist, dürfen wir die
Gruppengesetze an beliebigen Repräsentanten überprüfen.
Seien \(z=[a,b]\), \(w=[c,d]\) und \(t=[u,v]\).

\medskip\noindent\textbf{Assoziativität und Kommutativität.}
Die Assoziativität in \(M\) liefert
\[
  \begin{aligned}
    (z\oplus w)\oplus t
      &=[(a+c)+u,(b+d)+v]\\
      &=[a+(c+u),b+(d+v)]\\
      &=z\oplus(w\oplus t).
  \end{aligned}
\]
Hier sind bereits die jeweiligen Paare gleich. Ebenso folgt aus
der Kommutativität in \(M\)
\[
  z\oplus w=[a+c,b+d]=[c+a,d+b]=w\oplus z.
\]
Da jede Klasse eine solche Darstellung besitzt, gelten beide
Gesetze auf ganz \(K(M)\).

\medskip\noindent\textbf{Neutrales Element.}
Setze \(0_K:=[0,0]\). Dann gilt
\[
  [a,b]\oplus0_K=[a+0,b+0]=[a,b].
\]
Wegen der Kommutativität ist \(0_K\) auch auf der anderen
Seite neutral. Jedes Paar mit zwei gleichen Einträgen stellt
diese Null dar: \([r,r]=[0,0]\), denn \(r+0=r+0\).

\medskip\noindent\textbf{Inverse durch Vertauschen.}
Zum Paar \((a,b)\) betrachten wir \((b,a)\). Es ergibt sich
\[
  [a,b]\oplus[b,a]=[a+b,b+a]=[0,0]=0_K.
\]
Die mittlere Klassengleichheit folgt aus \(a+b=b+a\) und
der Kreuzsummenbedingung. Die Summe in umgekehrter Reihenfolge
hat wegen der Kommutativität denselben Wert. Jede Klasse besitzt
also ein beidseitiges Inverses.

Auch das Vertauschen ist unabhängig vom Repräsentanten. Aus
\([a,b]=[c,d]\) folgt \(a+d=b+c\), also \(b+c=a+d\).
Die letzte Gleichung besagt gerade \([b,a]=[d,c]\). Deshalb
dürfen wir die Negation auf Klassen definieren durch
\[
  -[a,b]:=[b,a].
\]

Damit erfüllt \((K(M),\oplus,0_K)\) alle Gesetze einer
abelschen Gruppe. Wir haben weder schon vorhandene Inverse in
\(M\) angenommen noch eine Subtraktion auf \(M\) verwendet.
Die neuen Inversen sind durch das Vertauschen der beiden
Darstellungskomponenten entstanden. Dieser Teil der Konstruktion
schließt mit \FormulaRefAuto{GroupCompletionAbelianGroup}[thm].

\par\Needspace{10\baselineskip}
\hypertarget{differences.reading.embedding}{}
\section*{Das ursprüngliche Monoid in der Gruppe}

Es bleibt zu zeigen, dass der neue Bereich die ursprünglichen
Elemente und ihre Operation bewahrt. Definiere
\[
  \iota\colon M\longrightarrow K(M),
  \qquad \iota(a):=[a,0].
\]
Jedem Element wird eine eindeutig bestimmte Klasse zugeordnet.
Die Formel definiert daher eine Funktion.

\medskip\noindent\textbf{Injektivität.}
Aus \(\iota(a)=\iota(b)\) folgt \([a,0]=[b,0]\).
Die Klassencharakterisierung liefert \(a+0=0+b\), also
\(a=b\). Verschiedene ursprüngliche Elemente bleiben somit
unterscheidbar.

\medskip\noindent\textbf{Erhaltung der Operation.}
Für \(a,b\in M\) gilt
\[
  \begin{aligned}
    \iota(a)\oplus\iota(b)
      &=[a,0]\oplus[b,0]\\
      &=[a+b,0+0]\\
      &=[a+b,0]=\iota(a+b).
  \end{aligned}
\]
Außerdem ist \(\iota(0)=[0,0]=0_K\). Damit ist \(\iota\)
eine injektive Abbildung, die Operation und neutrales Element
erhält. Der Einbettungssatz ist bewiesen.

\DifferencesEmbeddingDiagram

Genau genommen sind die Elemente von \(K(M)\) Mengen von
Paaren. Die ursprünglichen Elemente von \(M\) werden durch
ihr Bild \(\iota[M]\subseteq K(M)\) vertreten. Die
Abbildung \(\iota\) ist eine Bijektion von \(M\) auf dieses
Bild und erhält dort die gegebene Operation. In diesem präzisen
Sinn enthält die Gruppe eine Kopie des Monoids.

\medskip\noindent\textbf{Formale Differenzen werden wirkliche Differenzen.}
In der Gruppe definieren wir für \(z,w\in K(M)\)
\[
  w-z:=w\oplus(-z).
\]
Dann ist die anfangs nur beabsichtigte Lesart der Paare eine
bewiesene Gleichung:
\[
  \iota(a)-\iota(b)
    =[a,0]\oplus[0,b]=[a,b].
\]
Jedes Element der konstruierten Gruppe ist also eine Differenz
zweier eingebetteter Elemente von \(M\).

Nun ist auch jede Gleichung \(x\oplus z=w\) eindeutig
lösbar. Das Element \(w\oplus(-z)\) ist eine Lösung, denn
\[
  (w\oplus(-z))\oplus z
    =w\oplus((-z)\oplus z)=w.
\]
Erfüllt umgekehrt \(x\) die Gleichung \(x\oplus z=w\), so folgt
\[
  x=x\oplus0_K
    =x\oplus(z\oplus(-z))
    =(x\oplus z)\oplus(-z)=w\oplus(-z).
\]
Die Subtraktion löst damit die Aufgabe, die am Anfang der
Konstruktion stand.

Für \(M=\mathbb N\) erhalten wir insbesondere
\[
  [2,5]\oplus\iota(5)=[7,5]=[2,0]=\iota(2).
\]
Wegen \([2,5]=[0,3]=-\iota(3)\) ist dies die Antwort
\(x=-3\) auf die Ausgangsgleichung. Multiplikation und Ordnung
werden für diesen Schluss nicht benötigt.

\par\Needspace{10\baselineskip}
\hypertarget{differences.reading.product}{}
\section*{Zwei unabhängige Richtungen}

Betrachte \(M=\mathbb N^2\) mit der koordinatenweisen Operation
\[
  (a_1,a_2)+(c_1,c_2):=(a_1+c_1,a_2+c_2)
\]
und dem neutralen Element \((0,0)\). Assoziativität,
Kommutativität und Neutralität folgen in jeder Koordinate aus
den natürlichen Rechengesetzen. Ebenso ist dieses Monoid kürzbar:
Aus \((a_1,a_2)+(u_1,u_2)=(c_1,c_2)+(u_1,u_2)\) folgt
durch Kürzen in beiden Koordinaten \(a_1=c_1\) und \(a_2=c_2\).
Der Einbettungssatz ist daher anwendbar.

Ein Repräsentant in der neuen Konstruktion ist jetzt ein Paar
von Paaren: \((a,b)\), wobei \(a=(a_1,a_2)\) und
\(b=(b_1,b_2)\) natürliche Koordinaten haben. Für
\(c=(c_1,c_2)\) und \(d=(d_1,d_2)\) lautet die Relation
\[
  (a,b)\sim(c,d)
  \quad\Longleftrightarrow\quad
  \begin{cases}
    a_1+d_1=b_1+c_1,\\
    a_2+d_2=b_2+c_2.
  \end{cases}
\]
Die beiden Koordinaten werden unabhängig voneinander zu
Differenzen erweitert.

Das lässt sich als genaue Identifikation formulieren. Bezeichnen
wir die Klassen im eindimensionalen Modell \(K(\mathbb N)\)
mit \([r,s]_{\mathbb N}\), so definiert
\[
  \begin{aligned}
    F\colon K(\mathbb N^2)&\longrightarrow
                    K(\mathbb N)\times K(\mathbb N),\\
    F([a,b])&:=([a_1,b_1]_{\mathbb N},[a_2,b_2]_{\mathbb N})
  \end{aligned}
\]
eine wohldefinierte injektive Funktion: Zwei Eingabepaare
erfüllen genau dann die Relation in \(\mathbb N^2\), wenn
ihre beiden Koordinaten dieselben Klassen in \(K(\mathbb N)\)
liefern. Auch ist \(F\) surjektiv. Zu beliebigen Klassen
\([r,s]_{\mathbb N}\) und \([u,v]_{\mathbb N}\) wählen
wir \(a=(r,u)\) und \(b=(s,v)\); dann hat \([a,b]\)
das gewünschte Bild.

Schließlich erhält \(F\) die Addition: Die erste Koordinate
von \(F([a,b]\oplus[c,d])\) ist
\([a_1+c_1,b_1+d_1]_{\mathbb N}\), also die Summe der
ersten Bildkoordinaten; dieselbe Rechnung gilt für die zweite.
Die Nullklasse wird auf das Paar der Nullklassen abgebildet.
Wir erhalten somit eine bijektive strukturerhaltende Abbildung,
also einen Gruppenisomorphismus
\[
  K(\mathbb N^2)\cong K(\mathbb N)\times K(\mathbb N).
\]
Der formale Nachweis dieser Abbildung steht in
\FormulaRefAuto{FDProductNaturalSquareIsomorphism}[thm].
Das eindimensionale Modell lässt sich durch
\([r,s]_{\mathbb N}\mapsto[r,s]_{\mathbb Z}\) mit dem
Ganzzahlmodell aus Band~17 identifizieren: Beide Klassenkriterien
vergleichen dieselben natürlichen Kreuzsummen, und beide Additionen
addieren die Repräsentanten komponentenweise. In beiden Koordinaten
ergibt diese Identifikation
\[
  K(\mathbb N^2)\cong\mathbb Z^2.
\]
Die vollständige Identifikation mit dem Modell aus Band~17 wird in
\FormulaRefAuto{FDProductIntegerFinalIsomorphism}[thm] bewiesen.
Beispielsweise stellt das Paar
\(((2,5),(5,1))\) den Vektor \((-3,4)\) dar.
Die allgemeine Konstruktion ergänzt hier zwei unabhängig
veränderliche Größen um ihre additiven Gegenstücke.

\par\Needspace{10\baselineskip}
\hypertarget{differences.reading.assumptions}{}
\section*{Warum die Voraussetzungen genau passen}

Die Rolle der einzelnen Gesetze ist im Beweis sichtbar.
Assoziativität und Kommutativität erlauben das Umordnen der
Kreuzsummen; Kommutativität sichert bereits die Reflexivität
der gewählten Relation. Das neutrale Element liefert die
Nullklasse und die Einbettung \(a\mapsto[a,0]\).
Kürzbarkeit wird entscheidend bei der Transitivität gebraucht.
Der spätere Injektivitätsbeweis benutzt dann die bereits
bewiesene Klassencharakterisierung.

Ohne Kürzbarkeit kann diese einfache Relation tatsächlich
scheitern. Auf \(M=\{0,1\}\) sei die Operation das Maximum
zweier Elemente. Sie ist assoziativ und kommutativ und hat
das neutrale Element \(0\), ist aber nicht kürzbar:
\[
  \max(0,1)=\max(1,1),\qquad 0\ne1.
\]
Für die zugehörige Kreuzsummenrelation gelten
\[
  (0,0)\sim(1,1),\qquad (1,1)\sim(0,1),
  \qquad (0,0)\not\sim(0,1).
\]
Die erste Beziehung vergleicht zweimal \(\max(0,1)\), die
zweite \(\max(1,1)\) mit \(\max(1,0)\). In der dritten
wären dagegen \(\max(0,1)=1\) und \(\max(0,0)=0\)
gleichzusetzen. Die Transitivität ist also verletzt.
Die Tabellen weisen zuerst die Monoidgesetze nach
(\FormulaRefAuto{FDMaxCommutativeMonoid}[thm]) und führen dann
den Gegenbeweis zur Transitivität aus
(\FormulaRefAuto{FDMaxNotTransitive}[thm]).

Kürzbarkeit ist außerdem für jede Einbettung in eine Gruppe
notwendig. Angenommen, eine injektive Abbildung \(j\) erhält
die Operation und es gilt \(a+c=b+c\). Dann sind in der
Zielgruppe \(j(a)+j(c)\) und \(j(b)+j(c)\) gleich.
Durch Addition des Inversen von \(j(c)\) folgt
\(j(a)=j(b)\); Injektivität ergibt \(a=b\).
Für eine Einbettung in eine \emph{abelsche} Gruppe ist auch
Kommutativität notwendig: Dort gilt
\[
  j(a+b)=j(a)+j(b)=j(b)+j(a)=j(b+a),
\]
und die Injektivität liefert \(a+b=b+a\).
Die Voraussetzungen beschreiben damit genau die Monoide,
für die eine Einbettung in eine abelsche Gruppe möglich ist.
Die beiden notwendigen Bedingungen werden in
\FormulaRefAuto{FDNecessaryConditions}[thm] zusammengeführt;
die abschließende Äquivalenz steht in
\FormulaRefAuto{FDNecessaryEmbeddingCharacterization}[thm].

\par\Needspace{18\baselineskip}
\medskip\noindent\textbf{Anschluss an die beiden Bände.}
Der \DifferencesMainLink{} und die zugehörigen
\DifferencesProofsLink{} verwenden für die ursprüngliche
Monoidoperation das Zeichen \(\star\) und für ihr neutrales
Element \(e\). Unser allgemeines \(+\) und \(0\) sind
eine andere Schreibweise für dieselben Daten.

Im Zahlenfall liefert dieser Satz die additive Grundlage.
Die \IntegersReadingLink{} zu Band~17 entwickelt daraus
den konkreten Zahlbereich mit weiteren Operationen und seiner
Ordnung. Die dortigen \IntegersProofsLink{} führen den
formalen Modellnachweis bis zur Ganzzahlaxiomatik aus.
Der allgemeine Satz hier erklärt, welcher Teil dieses Aufbaus
allein aus den Monoidgesetzen und der Kürzbarkeit folgt:
eine abelsche Gruppe, eine Einbettung und die Darstellung
jedes neuen Elements als Differenz zweier alter Elemente.
